\section{Random Thoughts at Fall Harvest 2012}

\paragraph{Other Web Sites}

Recently a friend of Gornahoor sent me a link to another tradition-oriented web site. That is fine and there should be 10,000 of them. Nevertheless, sharing links is a form of intellectual laziness as specificity is what counts. I would add that accumulating quotes or proof texts falls in the same category. What I would prefer is a link to a specific text, with some context, including what is good or bad about it.

In general, however, there is little purpose in railing against specific elements in the modern world, as the problem is systemic. We know by now the factors that led up to it and that only a spiritual and intellectual conversion is what counts. That is where the focus needs to be.

Meanwhile, there are certainly Signposts as guides to determine if a given site or system of thought is leading in the right direction: viz., does it

\begin{enumerate}
\item know what Tradition is?
\item know what the Right is?
\item know what normal is?
\end{enumerate}

\paragraph{Maurras and the New Right}

I have postponed the review of Dugin's book until I have time to reconsider his approach. Some commentators have pointed out that adapting concepts from post-modern writers does not \emph{ipso facto} disqualify his argument.

To clarify the point of translating Maurras: the modern mind regards reason and freedom as revolutionary forces that have, and will continue to, transform man and society. Maurras is the antidote in that, for him, reason and freedom confirm the traditional order of things. The modern mind believes that by transforming social structure, man himself will change. Maurras, however, in looking just at the facts, free from the ``fantasies of the heart, wishes of the imagination, the needs, the amenities and personal interests'', knows man as he is, not as the revolutionary imagines him to be.

In other words, Maurras knows the Logos, at least in its appearance if not in it transcendence. Nevertheless, Benoist rejects Maurras, whom he had followed in his youth. Specifically, Benoist in his own words rejects Logos in favor of Mythos. I don't think many of his followers grasp that, believing instead that Benoist is making an intellectual argument rather than framing a myth. I don't believe that myth. Keep in mind that a logically possibility of manifestation will actually manifest. I don't see the myth he proposes as logically possible, although I welcome discussion.

\paragraph{Audrey}

Señorita Chicalinda's contributions\footnote{\url{https://gornahoor.net/?author=158}} have not been well received, although they are much more widely read than the Maurras translation. We will have to reconsider her future contributions. Nevertheless, as a perspective, it cannot be dismissed \emph{tout court}. I know Audrey, and she does have an unusual, but not unintelligent, mind. She not only thinks ``out of the box'', I doubt she even knows where the box is. I've been thinking about her comment, ``I also live outside of the world of library, university and home.''

On the one hand, they stand for learning and stability. On the other, the library is a testament to endless discussion without resolution, the university at best a place to catalog knowledge rather than teach it, and the home, a place for emotional abuse or worse. They are \emph{tulpas}, thought forms created out of human imagination, that then overpower their creators.

Outside of library, university and home, even outside the natural state of the lion, she sees something beyond nature, a sexuality that is neither pleasurable nor procreative. She swears an ``oath not to use her body as a vehicle of another's desires.'' As Gurdjieff says, ``men use women as a Kleenex.'' Nevertheless, we see contracepted women claiming that as a civil right. Even Evola saw women as mere receptacles for the band of warriors. In this area, Miguel Serrano\footnote{\url{https://gornahoor.net/?p=70}} seems to have understood a little more:

\begin{quotex}
Therefore, my love, do not take possession of my body. Let us not create children of the flesh. I will make you pregnant with the son of death. And we will both remain virgins.
\end{quotex}



\flrightit{Posted on 2012-10-03 by Cologero}

\begin{center}* * *\end{center}

\begin{footnotesize}
\begin{sffamily}
\texttt{Dominion on 2012-10-04 at 22:07 said:}

The comment that ''a logical possibility of manifestation will actually manifest'' brings immediately to mind the theory that all possible universes exist... the findings of science and Tradition and their similarity never cease to amaze and it is apparent why a positivist like Comte can be so influential to this site.

I wonder at the relationship between women and men described here. Recognizing that the relationship of ''pleasure and procreation'' is a relationship of egos, impermanent and changing, what would you say is the place of such relationships in the life of the person seeking to live by Tradition? Alain Danielou in his `Virtue, Success, Pleasure and Liberation' describes the tradition of the householder to marry and have children in his youth, and pleasure and procreation certainly played a role there. ''For everything, there is a season... ''


\hfill

\end{sffamily}
\end{footnotesize}

